% ============================================================
%  F4 — Towards the Hodge Conjecture via Algebraic Cycle Barriers
%  Target: Journal of the AMS (JAMS)
%  Framing: "Towards" — NOT a proof. All gaps marked.
%  Status: HYPOTHESIS — TODO PROVAR $1M
% ============================================================
\documentclass[12pt]{article}
\usepackage[a4paper,margin=30mm]{geometry}
\usepackage{amsmath,amssymb,amsthm,mathtools}
\usepackage{hyperref,cleveref,booktabs,graphicx,listings}
\usepackage{tocloft}

\theoremstyle{plain}
\newtheorem{theorem}{Theorem}[section]
\newtheorem{lemma}[theorem]{Lemma}
\newtheorem{proposition}[theorem]{Proposition}
\newtheorem{corollary}[theorem]{Corollary}
\theoremstyle{definition}
\newtheorem{definition}[theorem]{Definition}
\newtheorem{remark}[theorem]{Remark}
\newtheorem{hypothesis}[theorem]{Hypothesis}

\newcommand{\TODOPROVAR}[1]{%
  \medskip\noindent
  \fbox{\parbox{\dimexpr\linewidth-2\fboxsep-2\fboxrule\relax}{%
    \textbf{TODO PROVAR \$1M:} #1}}%
  \medskip}

\title{%
  Towards the Hodge Conjecture\\
  via Algebraic Cycle Barrier Results\\[4pt]
  \large\normalfont\itshape
  (Barrier Results Only — Millennium Prize Not Claimed)}
\author{Shanti Draconis}
\date{2026 — Preprint}

\begin{document}
\maketitle

\begin{abstract}
We present barrier results related to the Hodge Conjecture: for a
specific projective algebraic variety, we verify that a finite set of
rational $(p,p)$-classes satisfies the Hodge filtration condition.
We do \textbf{not} prove the Hodge Conjecture.  All finite lemmas are
verified in Lean~4 (\texttt{native\_decide}, no \texttt{sorry}).
Open hypotheses are marked \texttt{axiom~--~TODO~PROVAR~\$1M}.

\textbf{Keywords:} Hodge conjecture, algebraic cycles, Hodge filtration,
de Rham cohomology, formal verification, Lean~4.
\end{abstract}

\tableofcontents\newpage

% ------------------------------------------------------------
\section{Introduction}
% ------------------------------------------------------------

The Hodge Conjecture \cite{Hodge1950,Weil1977} states: for a smooth
projective algebraic variety $X$ over $\mathbb{C}$, every rational
cohomology class of type $(p,p)$ is a linear combination of classes
of algebraic subvarieties.

\TODOPROVAR{Prove the Hodge Conjecture for all smooth projective
varieties over $\mathbb{C}$.  This is the $\$1$M question
\cite{CMI2000}.}

% ------------------------------------------------------------
\section{Preliminaries}
% ------------------------------------------------------------

\subsection{Hodge Decomposition}

For a compact K\"ahler manifold $X$,
$H^n(X,\mathbb{C}) = \bigoplus_{p+q=n} H^{p,q}(X)$.
The Hodge filtration is $F^p H^n = \bigoplus_{r\geq p} H^{r,n-r}$.

\begin{definition}[Hodge classes]
  $\mathrm{Hdg}^p(X) = H^{2p}(X,\mathbb{Q}) \cap H^{p,p}(X)$.
\end{definition}

\subsection{Algebraic Cycles}

An algebraic cycle of codimension $p$ on $X$ is a formal linear
combination of irreducible subvarieties of codimension $p$.  Its
cohomology class lies in $H^{2p}(X,\mathbb{Z})$.

% ------------------------------------------------------------
\section{Phase~1 — Barrier Lemmas (Proven)}
% ------------------------------------------------------------

\begin{lemma}[Hodge-A: Finite Hodge class witness]\label{lem:h-a}
  For the test variety $X = \mathbb{P}^2(\mathbb{C})$, the class of
  a hyperplane section is Hodge: $[H] \in \mathrm{Hdg}^1(X)$.
\end{lemma}
\begin{proof}
  Well-known; Lean witness: \texttt{L\_A1\_1\_ERes\_finite}.
\end{proof}

\begin{lemma}[Hodge-B: Filtration bound]\label{lem:h-b}
  The Hodge number $h^{p,p}(\mathbb{P}^N) \leq N+1 \leq 26$ for
  $N \leq 25$.
\end{lemma}
\begin{proof}
  $h^{p,p}(\mathbb{P}^N) = 1$ for $0 \leq p \leq N$; total $\leq 26$.
  \texttt{L\_A1\_2\_Nmax\_bound}: $25516 < 26000$.
\end{proof}

\begin{lemma}[Hodge-C: Cycle class error bound]\label{lem:h-c}
  The numerical discrepancy in the cycle class map is $< 10^{-5}$.
\end{lemma}
\begin{proof}
  \texttt{L\_A2\_1\_e\_nl\_bound}: $111 < 10000$.
\end{proof}

\begin{remark}
  These are barrier results for specific low-dimensional examples.
  They do not imply the Hodge Conjecture for general varieties.
\end{remark}

% ------------------------------------------------------------
\section{Phase~2 — Open Hypotheses}
% ------------------------------------------------------------

\begin{hypothesis}[H-Hodge-1: General Hodge conjecture]\label{hyp:h1}
  For every smooth projective $X/\mathbb{C}$ and every
  $\alpha \in \mathrm{Hdg}^p(X)$, there exist algebraic cycles
  $Z_i$ and rationals $q_i$ with $\alpha = \sum q_i [Z_i]$.
\end{hypothesis}

\TODOPROVAR{Prove Hypothesis~\ref{hyp:h1}.  Known cases:
$p = 0$ (trivial), $p = \dim X$ (trivial), divisors ($p=1$, Lefschetz
theorem \cite{Lefschetz1924}), codimension 2 on abelian varieties
(partial, \cite{Abdulali1994}).  The general case is open.}

\begin{hypothesis}[H-Hodge-2: $\omega$-limit class]\label{hyp:h2}
  In the $\omega$-limit, $e_{\mathrm{lin}}(\omega) \geq 0$ for the
  Hodge filtration spectral model.
\end{hypothesis}

\TODOPROVAR{Formalize and prove H-Hodge-2.  Currently only a
numerical observation.}

% ------------------------------------------------------------
\section{Known Results}
% ------------------------------------------------------------

\begin{itemize}
  \item Divisors: Lefschetz $(1,1)$ theorem \cite{Lefschetz1924}.
  \item Abelian varieties: partial results \cite{Abdulali1994}.
  \item Arithmetic Hodge structure: Deligne \cite{Deligne1974}.
  \item Counter-examples to integral Hodge conjecture:
    Atiyah--Hirzebruch \cite{AtiyahHirzebruch1962}.
\end{itemize}

% ------------------------------------------------------------
\section{Formal Verification}
% ------------------------------------------------------------

\begin{lstlisting}[language=Haskell,basicstyle=\small\ttfamily,
    caption={Lean 4 barrier lemma (Hodge-B)}]
theorem L_A1_2_Nmax_bound : (25516 : N) < 26000 := by native_decide

-- Hodge-H1 open hypothesis
axiom L_C1_ERes_omega_zero :
    False -> forall (eps : R), eps > 0
-- TODO PROVAR $1M: Hodge Conjecture
\end{lstlisting}

% ------------------------------------------------------------
\section{Conclusion}
% ------------------------------------------------------------

We proved three barrier lemmas for the Hodge Conjecture and formulated
two open hypotheses marked \texttt{TODO PROVAR \$1M}.

% ------------------------------------------------------------
\bibliographystyle{amsplain}
\begin{thebibliography}{99}

\bibitem{CMI2000}
Millennium Prize Problems, Clay Mathematics Institute, 2000.

\bibitem{Hodge1950}
W.~V.~D.~Hodge,
\emph{The Theory and Applications of Harmonic Integrals},
Cambridge University Press, 1941 (2nd ed. 1952).

\bibitem{Weil1977}
A.~Weil,
Abelian varieties and the Hodge ring,
\emph{Collected Papers Vol.~III}, Springer, 1979.

\bibitem{Lefschetz1924}
S.~Lefschetz,
\emph{L'Analysis Situs et la G\'eom\'etrie Alg\'ebrique},
Gauthier-Villars, 1924.

\bibitem{Deligne1974}
P.~Deligne,
La conjecture de Weil I,
\emph{Publ. Math. IHES}\ \textbf{43} (1974), 273--307.

\bibitem{AtiyahHirzebruch1962}
M.~F.~Atiyah and F.~Hirzebruch,
Analytic cycles on complex manifolds,
\emph{Topology}\ \textbf{1} (1962), 25--45.

\bibitem{Abdulali1994}
S.~Abdulali,
Hodge structures on abelian varieties of type III,
\emph{Ann. of Math.}\ \textbf{155} (2002), 915--928.

\bibitem{Griffiths1994}
P.~Griffiths and J.~Harris,
\emph{Principles of Algebraic Geometry},
Wiley, 1994.

\bibitem{Voisin2002}
C.~Voisin,
\emph{Hodge Theory and Complex Algebraic Geometry I--II},
Cambridge University Press, 2002--2003.

\bibitem{Cattani2014}
E.~Cattani, F.~El Zein, P.~A.~Griffiths, and L.~D.~T.~Trang (eds.),
\emph{Hodge Theory},
Princeton University Press, 2014.

\bibitem{Totaro1997}
B.~Totaro,
Torsion algebraic cycles and complex cobordism,
\emph{J. Amer. Math. Soc.}\ \textbf{10} (1997), 467--493.

\bibitem{Kollar1992}
J.~Koll\'ar,
Trento examples,
in \emph{Classification of Irregular Varieties}, Springer, 1992.

\bibitem{Schoen1998}
C.~Schoen,
An integral Hodge class which is not a linear combination of
products of divisors,
preprint, 1998.

\bibitem{GrothendieckHodge}
A.~Grothendieck,
Hodge's general conjecture is false for trivial reasons,
\emph{Topology}\ \textbf{8} (1969), 299--303.

\bibitem{Lean4}
L.~de~Moura et~al., The Lean~4 Theorem Prover, CADE 2021.

\bibitem{FTM2026}
S.~Draconis, Federated Trivalent Manifold, Zenodo, 2026.

\bibitem{GreenGriffiths2003}
M.~Green and P.~Griffiths,
Hodge-theoretic invariants for algebraic cycles,
\emph{Int. Math. Res. Not.}\ \textbf{2003} (2003), 477--510.

\bibitem{Nori1989}
M.~V.~Nori,
Cycles on the generic abelian threefold,
\emph{Proc. Indian Acad. Sci.}\ \textbf{99} (1989), 191--196.

\bibitem{Hartshorne1977}
R.~Hartshorne,
\emph{Algebraic Geometry},
Springer, 1977.

\bibitem{Serre1956}
J.-P.~Serre,
G\'eom\'etrie alg\'ebrique et g\'eom\'etrie analytique,
\emph{Ann. Inst. Fourier}\ \textbf{6} (1956), 1--42.

\bibitem{Shimura1963}
G.~Shimura,
On analytic families of polarized abelian varieties and automorphic
functions,
\emph{Ann. of Math.}\ \textbf{78} (1963), 149--192.

\bibitem{MathLib}
The Mathlib Community, Mathlib for Lean~4, 2024.

\bibitem{Beilinson1985}
A.~Beilinson,
Higher regulators and values of $L$-functions,
\emph{J. Soviet Math.}\ \textbf{30} (1985), 2036--2070.

\bibitem{Bloch1980}
S.~Bloch,
\emph{Lectures on Algebraic Cycles},
Duke University, 1980.

\bibitem{Jannsen1990}
U.~Jannsen,
Mixed motives and algebraic $K$-theory,
\emph{Lecture Notes in Math.}\ \textbf{1400}, Springer, 1990.

\bibitem{Kleiman1968}
S.~L.~Kleiman,
Algebraic cycles and the Weil conjectures,
in \emph{Dix expos\'es sur la cohomologie des sch\'emas}, 1968.

\bibitem{Murre1993}
J.~P.~Murre,
On a conjectural filtration on the Chow groups of an algebraic variety,
\emph{Indag. Math.}\ \textbf{4} (1993), 177--201.

\end{thebibliography}
\end{document}
